\documentclass{scrartcl}
\usepackage{graphicx}  % required for inserting images

% for cyrillic and Bulgarian language
\usepackage[utf8]{inputenc}
\usepackage[T2A]{fontenc}
\usepackage[bulgarian]{babel}

% math packages
\usepackage{amsmath}
\usepackage{float}
\usepackage{amsfonts}
\usepackage{amssymb}
\usepackage{amsthm}
\usepackage{cancel}

\usepackage{ragged2e}  % adds FlushLeft
\usepackage{booktabs}
\usepackage{tabularx}
\usepackage{multirow}   % for complex tables

\usepackage{listings}
\usepackage{xcolor}

\definecolor{codegreen}{rgb}{0,0.6,0}
\definecolor{codegray}{rgb}{0.5,0.5,0.5}
\definecolor{codepurple}{rgb}{0.58,0,0.82}
\definecolor{backcolour}{rgb}{0.95,0.95,0.92}

\usepackage[margin=3cm]{geometry}

\lstdefinestyle{mystyle}{
    %caption=Example C++,
    label={lst:listing-cpp},
    language=C++,
    backgroundcolor=\color{backcolour},   
    commentstyle=\color{codegreen},
    keywordstyle=\color{magenta},
    numberstyle=\tiny\color{codegray},
    stringstyle=\color{codepurple},
    basicstyle=\ttfamily\footnotesize,
    breakatwhitespace=false,         
    breaklines=true,                 
    keepspaces=true,                 
    numbers=left,       
    numbersep=5pt,                  
    showspaces=false,                
    showstringspaces=false,
    showtabs=false,                  
    tabsize=2,
}
\lstset{style=mystyle}

\usepackage[colorlinks=true, linkcolor=blue, urlcolor=cyan]{hyperref}

\newtheorem{theorem}{Теорема}
\newtheorem{definition}{Дефиниция}

\title{Модели и методи за проектиране на потребителски интерфейс}
\author{Ивайло Андреев + DeepSeek-V3}
\date{13 юни 2026}

\graphicspath{{./}}

\begin{document}
\maketitle
\tableofcontents
\newpage

\begin{FlushLeft}

\section{Основни модели и методи при създаване на потребителски интерфейс}

Дейността по проектиране на човеко-машинен интерфейс включва проектиране, създаване и оценяване на интерактивни компютърни системи, предназначени за използване от хора. Основно значение има взаимодействието между един или повече потребители с една или повече програми.

\subsection{Подходи}

Съществуват четири основни подхода в проектирането на потребителски интерфейси:
\begin{enumerate}
    \item \textbf{Потребителят в центъра (User-centred design – UCD)} – потребителят е основен източник на изисквания, а проектантът ги превежда на езика на интерфейса. Този подход поставя потребителя във фокуса на целия процес на разработка.
    \item \textbf{Действията в центъра (Activity-centred design – ACD)} – насочен е не към глобалните цели, а към конкретните стъпки при всяко действие. Анализира се какво точно прави потребителят, а не какво иска да постигне.
    \item \textbf{Системно проектиране} – акцентът е върху крайната система като цяло, включително хардуер, софтуер и работна среда.
    \item \textbf{Артистично проектиране (Genius design)} – разчита се на опита, интуицията и творчеството на дизайнера, а ролята на потребителя е единствено да потвърди готовия дизайн.
\end{enumerate}

\subsection{Основни процеси}

Основните дейности в процеса на проектиране са четири, като те се изпълняват итеративно:
\begin{enumerate}
    \item \textbf{Определяне на нуждите и установяване на изискванията} – чрез анкети, интервюта, наблюдения и анализ на съществуващи системи.
    \item \textbf{Проектиране на алтернативни решения} – създаване на няколко възможни варианта за интерфейс.
    \item \textbf{Избор между алтернативите (оценяване)} – сравняване на вариантите по критерии като използваемост, ефективност и удовлетвореност на потребителите.
    \item \textbf{Създаване на прототип} – изграждане на работещ модел на интерфейса за тестване и валидиране.
\end{enumerate}

\subsection{Анализ на задачите}

Анализът на задачите е систематичен метод за изучаване на човешките дейности. Той помага да се разбере как потребителите извършват дадена работа и какви са техните потребности.

Основните методи са:
\begin{itemize}
    \item \textbf{Task Analysis (TA)} – простите действия се разглеждат като „черни кутии“. Анализират се входовете и изходите на всяко действие без детайлно разкриване на вътрешните процеси.
    \item \textbf{Cognitive Task Analysis (CTA)} – анализира задачи, които изискват значителна умствена дейност: решаване на проблеми, вземане на решения, внимание, работна памет и преценка. Целта е да се представят когнитивните стратегии, използвани от експертите.
    \item \textbf{Hierarchical Task Analysis (HTA)} – всички действия се представят като йерархия, където сложните задачи се разлагат на по-прости подзадачи. За всяка подзадача се определят условията за успешно изпълнение.
\end{itemize}

\subsection{Специфициране на взаимодействията}

Взаимодействието между потребителя и интерфейса може да бъде специфицирано чрез различни формални и полуформални средства:
\begin{itemize}
    \item \textbf{Формални граматики} – BNF (Backus-Naur Form), многозначни граматики, които описват синтаксиса на потребителските команди.
    \item \textbf{Диаграми на преходите} – дървета на диалоговите кутии, графи на състоянията (крайни автомати) за визуализация на възможните потребителски пътеки.
    \item \textbf{User Action Notation (UAN)} – специализиран инструмент за специфициране на потребителски действия.
\end{itemize}

\textbf{User Action Notation (UAN)} е инструмент, който не зависи от конкретния потребител или технологията на интерфейса. Той цели да покаже връзката между действията и интерфейсните елементи, като описва интерфейсите прецизно, подробно и еднозначно. Основната градивна единица е действие на потребител в даден контекст. Всяко действие се описва чрез сценарии, спомагателни бележки и диаграма на преходите между състоянията при конкретни условия. Чрез UAN могат да се репрезентират интерфейси с директно манипулиране (Direct Manipulation Interfaces).

\subsection{Основни техники и инструментални средства}

Инструменталните средства за проектиране на интерфейси се категоризират както следва:
\begin{itemize}
    \item \textbf{Средства с общо предназначение} – PowerPoint, Visio, InDesign. Използват се за бързи скици и презентации на идеи.
    \item \textbf{Средства за спецификация на софтуерни системи} – базирани на UML (диаграми на класове, диаграми на дейности, диаграми на състоянията).
    \item \textbf{Специфични системи за интерфейси} – Axure RP, LucidChart, ProtoShare. Позволяват създаване на интерактивни wireframe-и и прототипи.
    \item \textbf{Системи за прототипиране} – Balsamiq Mockups, Figma, Adobe XD. Фокусирани са върху бързото създаване на кликваеми прототипи.
    \item \textbf{Визуални среди за програмиране} – Visual Studio, Qt Designer, Android Studio – позволяват изграждане на напълно функциониращи интерфейси.
\end{itemize}

Критерии за избор на подходящо средство са: поддръжка на шаблони (widgets), богатство на интерфейсните компоненти, подходяща софтуерна архитектура и удобство при управление на проекта.

\section{Проектиране на графичен интерфейс}

\subsection{Интерактивни стилове и техники}

Формите на взаимодействие на потребителя с потребителския интерфейс се класифицират в пет основни стила:
\begin{itemize}
    \item \textbf{Директно въвеждане (Direct Manipulation)} – потребителят работи директно с визуални обекти на екрана. Пример: за да изтрие файл, потребителят го влачи с мишката върху иконата на кошчето. Характеризира се с незабавна визуална обратна връзка.
    \item \textbf{Избор от меню (Menu Selection)} – потребителят избира команда от списък с възможности. Менютата могат да бъдат падащи, изскачащи, контекстни или хоризонтални.
    \item \textbf{Попълване на форми (Form Fill-in)} – потребителят попълва данни в предварително дефинирани полета, подобно на хартиени формуляри. Често се използва при въвеждане на структурирана информация.
    \item \textbf{Команден език (Command Language)} – потребителят въвежда текстови команди и параметри. Осигурява висока гъвкавост и скорост за опитни потребители, но изисква запомняне на синтаксис.
    \item \textbf{Естествен език (Natural Language)} – потребителят използва нормална реч (писмена или говорна) за комуникация със системата. Примери са чатботове, виртуални асистенти (Siri, Alexa) и системи за въпроси и отговори.
\end{itemize}

\subsection{Отчитане на психологичните особености на потребителите}

Отчитането на познавателните способности на потребителите е от съществено значение за успешното проектиране. Разграничават се две основни групи:
\begin{itemize}
    \item \textbf{Експериментални (експресивни) особености} – свързани са с начина на мислене, сравнение, вземане на решения и креативност. Те варират при различните потребители и зависят от контекста.
    \item \textbf{Рефлективни особености} – отнасят се до начина на действие и реакции при различни събития, както и до емоционалното въздействие на интерфейса. Рефлективният опит включва осмисляне, учене и културни ценности.
\end{itemize}

Проектантът трябва да взема предвид когнитивните ограничения на човека – ограничен капацитет на работната памет (7±2 единици), ограничено внимание и склонност към грешки. Интерфейсите трябва да поддържат както експерименталната, така и рефлективната познавателна дейност.

\subsection{Концептуални модели и метафори}

\textbf{Концептуалните модели} описват системата като множество от интегрирани идеи и концепции – какво трябва да прави, как да изглежда и какво поведение да има в конкретни ситуации. Те помагат на потребителя да предвиди поведението на системата и да разбере как да постигне целите си.

Видове концептуални модели:
\begin{itemize}
    \item \textbf{Базирани на дейностите} – описват основните дейности, които потребителят извършва (комуникация, навигация, редактиране, търсене).
    \item \textbf{Базирани на обектите} – фокусират се върху конкретни обекти, техните свойства и операции, които могат да се прилагат върху тях.
    \item \textbf{Смесени} – комбинация от горните два подхода.
\end{itemize}

\textbf{Метафорите} са мощен инструмент за изграждане на концептуални модели. Те пренасят познати аспекти от физическия свят в дигиталната среда. Примери: „работен плот“ с папки и документи, „кошче“ за изтрити файлове, „мащабиране“ като увеличение на изображение. Метафорите обогатяват знанието на потребителя, като му позволяват да използва предишен опит, но трябва да се прилагат внимателно – когато метафората е непоследователна, тя обърква, вместо да помага.

\subsection{Методи и средства за реализация}

За реализация на графични интерфейси се използват следните методи:

\begin{itemize}
    \item \textbf{Сценарии (Storyboarding)} – описват как системата реагира при конкретна последователност от стъпки в определена среда. Предимства:
    \begin{itemize}
        \item Фокусират се върху цялата система, а не върху отделни компоненти.
        \item Абстрахират се от конкретните детайли на интерфейса (менюта, бутони).
        \item Свързват всички роли в дадена дейност за постигане на обща цел.
    \end{itemize}
    \item \textbf{Хартиени прототипи} – интерактивен макет, направен от изрезки хартия (бутони, менюта, текстови полета). Човек („компютър“) симулира работата на системата, като размества елементите според действията на потребителя. Предимства:
    \begin{itemize}
        \item Бързи, лесни и евтини за създаване.
        \item Гъвкави – промените се правят на момента.
        \item Всички участници (дизайнери, програмисти, потребители) могат да участват.
        \item Фокусът е върху пълната система, а не върху визуалните детайли.
    \end{itemize}
    \item \textbf{Компютърни прототипи} – интерактивна софтуерна симулация с висока прецизност на визуалното представяне, но често с ограничена функционалност (хоризонтален прототип). При липса на реализирана сървърна част се използва техниката \textbf{Wizard of Oz} – скрит човек, който симулира отговорите на системата.
    \item \textbf{Смесени прототипи} – комбинация от хартиени и компютърни елементи. Например част от екраните са реализирани на компютър, а други се симулират на хартия.
\end{itemize}

\section{Разработка на използваем графичен интерфейс: техники базирани на експерименти}

Използваемостта (usability) е основна характеристика на качествения интерфейс. Според стандарта ISO 9241-11, използваемостта се определя от три критерия: ефективност (резултатност), ефикасност (експедитивност) и удовлетвореност.

\subsection{Критерии за използваемост}

По-детайлно, използваемостта включва:
\begin{itemize}
    \item \textbf{Резултатност} – доколко интерфейсът помага на потребителя да извърши необходимите дейности напълно и точно.
    \item \textbf{Експедитивност} – дейностите се извършват с минимални усилия и бързо.
    \item \textbf{Безопасност} – интерфейсът предпазва потребителя от грешки и му позволява лесно да се възстанови (напр. бутон „Отмяна“).
    \item \textbf{Полезност} – каква част от реалните дейности на потребителя се поддържат от интерфейса.
    \item \textbf{Леснота на изучаване} – нов потребител може бързо да се научи да работи с интерфейса.
    \item \textbf{Възприемчивост (запомняемост)} – след период на неизползване потребителят лесно си припомня как да работи със системата.
\end{itemize}

\subsection{Техники за оценка, базирани на експерименти}

Провеждат се експерименти с типични потребители („лабораторни плъхове“) в специално създадени условия или в естествената им работна среда.

Основни методи за събиране на данни:
\begin{itemize}
    \item \textbf{Въпросници} – събират субективно мнение на потребителите за интерфейса (напр. SUS – System Usability Scale).
    \item \textbf{Наблюдение} – директно или чрез видеозапис се наблюдава как потребителите работят със системата, отбелязват се трудностите и времето за изпълнение на задачите.
    \item \textbf{Видеозаписи} – позволяват многократен анализ и детайлно проучване на потребителското поведение.
    \item \textbf{Инжектиране на „шпионски“ код} – записване на кликвания, движения на мишката, натиснати клавиши и грешки. Този метод дава обективни количествени данни от голям брой потребители.
    \item \textbf{Think Aloud протокол} – потребителите разсъждават на глас, докато работят, описвайки своите мисли, намерения и съмнения. Това разкрива когнитивните процеси зад действията.
\end{itemize}

\subsection{Когнитивни модели за предсказване на производителността}

\textbf{GOMS} (Goals, Operators, Methods, and Selection rules) е семейство от модели за предсказване на човешкото поведение. Те се използват за идентифициране и елиминиране на неефикасни човешки дейности. Семейството включва четири основни модела:

\begin{table}[H]
\centering
\caption{Сравнение на моделите от семейството GOMS}
\begin{tabular}{|p{2.2cm}|p{2cm}|p{2.2cm}|p{2.2cm}|p{2.2cm}|}
\hline
\textbf{Характеристика} & \textbf{KLM} & \textbf{CMN-GOMS} & \textbf{NGOMSL} & \textbf{CPM-GOMS} \\
\hline
Архитектура & Проста & Model Human Processor & Cognitive Complexity Theory & Model Human Processor + опит \\
\hline
Йерархия на цели & Неявна & Явна & Неявна & Неявна \\
\hline
Моделира учене & Не & Не & Да & Не \\
\hline
Моделира паралелни процеси & Не & Не & Не & Да \\
\hline
Време за умствена дейност & Да (оператор M) & Не & Да & Да (малко) \\
\hline
Нотация & Примитивни оператори & Език за програмиране & Естествен език & PERT диаграми \\
\hline
\end{tabular}
\end{table}

\textbf{KLM-GOMS (Keystroke-Level Model)} предсказва времето за изпълнение на задачи чрез сумиране на времената на прости оператори:
\begin{itemize}
    \item \textbf{K} (Keystroke) – натискане на клавиш или бутон на мишка (0.28 сек).
    \item \textbf{P} (Pointing) – насочване на мишката към цел (1.1 сек).
    \item \textbf{H} (Home) – преместване на ръката между клавиатура и мишка (0.4 сек).
    \item \textbf{M} (Mental) – умствена подготовка за следващо действие (1.35 сек).
    \item \textbf{R} (Response) – време за отговор на системата (забавяне).
\end{itemize}

\subsection{Когнитивен обход (Cognitive Walkthrough)}

Това е метод за инспекция на използваемостта, фокусиран върху леснотата на изучаване. Екип от оценители (обикновено 3–5 души) симулира работата на потребител, който за първи път използва системата. За всяка задача се задават четири въпроса:
\begin{enumerate}
    \item Ще се опита ли потребителят да постигне правилния ефект?
    \item Ще забележи ли потребителят, че правилното действие е налице?
    \item Ще свърже ли потребителят правилното действие с целта си?
    \item Ако извърши правилното действие, ще получи ли потребителят обратна връзка, че напредва към целта?
\end{enumerate}
Резултатите се записват като проблеми и препоръки за корекция.

\section{Разработка на мултимедиен графичен интерфейс}

Мултимедийният интерфейс интегрира различни типове медии: цвят, звук, текст, графика и анимация. Всеки от тези елементи има свои закономерности на проектиране.

\subsection{Проектиране на цветове}

Цветът има множество функции в интерфейса:
\begin{itemize}
    \item Привлича вниманието към важни елементи.
    \item Подобрява навигацията и скоростта на търсене.
    \item Показва взаимовръзки (групиране на обекти).
    \item Дава обратна връзка (напр. червено за грешка, зелено за успех).
\end{itemize}

Основни цветови модели:
\begin{itemize}
    \item \textbf{RGB} (Red, Green, Blue) – адитивен модел за екрани. Всеки цвят се задава с интензитет от 0 до 255.
    \item \textbf{CMYK} (Cyan, Magenta, Yellow, Black) – субтрактивен модел за печат.
    \item \textbf{HSL} (Hue, Saturation, Lightness) – интуитивен модел: оттенък (0–360°), наситеност (0–100\%) и светлина (0–100\%).
\end{itemize}

\textbf{Препоръки за използване на цвят:}
\begin{itemize}
    \item Цветовете да се прилагат последователно – един и същ цвят винаги да означава едно и също.
    \item Да се използват най-много осем цвята със специално значение.
    \item Да се избягва червено и синьо за шрифтове, както и червено-синя комбинация поради стереоскопичен ефект.
    \item Да се осигури достатъчно голям контраст между цвета на символите и цвета на фона.
    \item Препоръчва се тъмен фон за текст, но бял фон дава добра четливост.
\end{itemize}

\textbf{Цветови комбинации (четливост):}

\begin{table}[H]
\centering
\caption{Четливост на комбинации фон – цвят на знаците}
\begin{tabular}{|c|c|c|c|c|c|c|c|c|}
\hline
\multirow{2}{*}{\textbf{Цвят на фона}} & \multicolumn{8}{c|}{\textbf{Цвят на знаците}} \\
\cline{2-9}
 & Черно & Бяло & Пурпурно & Синьо & Циан & Зелено & Жълто & Червено \\
\hline
Черно & – & + & – & + & + & + & + & – \\
\hline
Бяло & + & – & + & + & + & + & + & + \\
\hline
Пурпурно & – & + & – & – & + & – & + & – \\
\hline
Синьо & – & + & – & – & + & – & + & – \\
\hline
Циан & + & + & + & + & – & + & + & + \\
\hline
Зелено & – & + & – & – & + & – & + & – \\
\hline
Жълто & + & + & + & + & + & + & – & + \\
\hline
Червено & – & + & – & – & + & – & + & – \\
\hline
\end{tabular}
\end{table}
(+ означава добра четливост, – означава лоша четливост)

\textbf{Цветови схеми:}

\begin{figure}[H]
\centering
\includegraphics[width=0.9\textwidth]{pictures/color_schemes.png}
\caption{Основни цветови схеми: монохроматична, аналогична, комплементарна, split-комплементарна, триадна, тетрадна и квадратна}
\label{fig:colorschemes}
\end{figure}

\subsection{Проектиране на звуци}

Звукът се използва в три основни форми:
\begin{itemize}
    \item \textbf{Звукови ефекти} – сигнализират за събития (напр. звук при получено съобщение, сигнал за грешка).
    \item \textbf{Говор} – синтезирана или записана реч за навигация, инструкции или обратна връзка.
    \item \textbf{Музика} – създава емоционална атмосфера, използва се в игри, обучителни системи и видео.
\end{itemize}
Звукът придава емоционално въздействие, обогатява мултимедийните елементи и дава възможност на автора директно да комуникира с потребителя. Формати: MP3, WAV, OGG, AAC. При проектиране трябва да се осигури възможност за изключване или регулиране на звука, за да не се нарушава спокойствието на потребителя.

\subsection{Проектиране на текст}

Текстът е основното средство за предоставяне на подробна информация. Препоръки за проектиране на текст в интерфейса:
\begin{itemize}
    \item \textbf{Контраст} – висок контраст между текста и фона (напр. черен текст върху бял фон). Избягват се конфликтни цветови двойки.
    \item \textbf{Шрифт и размер} – размер под 10 пункта се чете трудно. За заглавия и кратки инструкции е подходящ ясен шрифт без сериф (Arial, Helvetica). За дълги параграфи – шрифт със сериф (Times, Georgia).
    \item \textbf{Формулиране} – използват се минимален брой думи, които предават точно смисъла. Избягват се съкращения; ако са неизбежни, се използват стандартни такива.
    \item \textbf{Параметри на знаците} (според ергономични изследвания):
    \begin{itemize}
        \item Височина на знаците: препоръчителна големина 3–5 mm при нормално разстояние на гледане.
        \item Дебелина на знаците: 10–20\% от височината.
        \item Ширина на знаците: 50–70\% от височината.
        \item Разстояние между знаците: минимум 10\% от височината.
    \end{itemize}
\end{itemize}

Текстови кодирания:
\begin{itemize}
    \item \textbf{ASCII} – 7-битово кодиране, представя латиница, цифри и основни символи.
    \item \textbf{UTF-8} – стандартно Unicode кодиране с променлива дължина (1 до 4 байта). Латинските букви заемат 1 байт, кирилицата – 2 байта. Поддържа всички съвременни писмености.
\end{itemize}

\subsection{Проектиране на графика}

Графиките (изображенията) представят визуална информация на малка площ. Двата основни типа са:

\begin{itemize}
    \item \textbf{Растерна графика} – изображението е двумерна матрица от цветови точки (пиксели). Формати: JPEG (снимки), PNG (прозрачност, без загуба), GIF (анимация, ограничени цветове). Предимство: детайлност. Недостатък: загуба на качество при мащабиране.
    \item \textbf{Векторна графика} – изображението се състои от геометрични примитиви (линии, криви, полигони), описани с математически формули. Формати: SVG, EPS, AI. Предимство: мащабиране без загуба на качество. Недостатък: трудно представяне на фотореалистични изображения.
\end{itemize}

При проектиране на графични елементи (икони, бутони, илюстрации) трябва да се осигури яснота, простота и последователност на стила.

\subsection{Проектиране на анимация}

Анимацията представя движещи се образи, които показват промяна на състоянието на обектите във времето. Приложения в интерфейса:
\begin{itemize}
    \item Визуализиране на процеси, които са трудни за обяснение статично (напр. зареждане, преход между екрани).
    \item Привличане на вниманието към промяна (напр. пулсиращ бутон).
    \item Подобряване на пространственото ориентиране (напр. анимирано преминаване от списък към детайли).
\end{itemize}
Формати: MP4, MOV, AVI, GIF, WebM. Принципи на добрата интерфейсна анимация: умереност (без прекалено много движение), плавност (достатъчна честота на кадрите – минимум 24 fps), целенасоченост (всяка анимация да носи смисъл, а не да е декоративна).

\section{Особености при създаване на интегриран интерфейс}

Интегрираният интерфейс съчетава различни стилове, медии и методи за взаимодействие в една цялостна система. Той трябва да осигурява безпроблемна работа на потребителя с разнородни компоненти.

\subsection{Методи за моделиране, насочени към крайния потребител}

Тези методи имат за цел да съберат истинските нужди и контекст на потребителите:
\begin{itemize}
    \item \textbf{Етнография} – продължително наблюдение на потребителите в тяхната естествена работна среда. Изследователят се потапя в ежедневните дейности, за да разбере неписаните правила и навици.
    \item \textbf{Метод на свързаността} – комбинира етнографския подход с инженеринга на изискванията. Създават се модели, които свързват социалните и техническите аспекти на работата.
    \item \textbf{Метод на контекстния дизайн} – систематичен подход, разработен от Хю Байер и Карен Холцблат. Включва четири фази: контекстно интервю, интерпретация на данните, създаване на модели (поток, последователност, културни модели), валидиране с потребителите.
    \item \textbf{Метод на участието (Participatory Design)} – потребителите са активни членове на дизайнерския екип от самото начало. Те не само дават обратна връзка, но и съвместно създават прототипи, участват в мозъчни атаки и вземат решения за функционалността.
\end{itemize}

\subsection{Екранен дизайн}

Екранният дизайн обхваща подреждането на всички визуални елементи – прозорци, менюта, икони, полета за текст, бутони. Ключови принципи (в духа на Едуард Тюфти):
\begin{itemize}
    \item Подпомагане на визуални сравнения.
    \item Показване на последователност.
    \item Показване на многомерни променливи.
    \item Интегриране на текст, графика и данни в един екран.
    \item Осигуряване на качество, уместност и цялостност на съдържанието.
    \item Пространствена близост на свързаните неща и непротиворечивост във времето.
\end{itemize}

\textbf{Изисквания към иконите:}
\begin{itemize}
    \item Ясно видими и лесно разграничаваеми от околните елементи.
    \item Всяка икона да има уникална форма или пиктограма, различима дори без цвят.
    \item Сходен стил (еднакъв размер, линия, детайлност) за всички икони в рамките на приложението.
    \item Избягване на излишни детайли – иконата трябва да комуникира само основното си значение.
\end{itemize}

\textbf{Промишлени стандарти (Styleguides):} Това са сборници от препоръки за проектиране на графични интерактивни повърхнини. Те описват визуалния израз на обектите (look) и тяхното управление (feel). Примери:
\begin{itemize}
    \item IBM SAA-CUA (Common User Access) – Basic Interface Design Guide и Advanced Interface Design Reference.
    \item Apple Human Interface Guidelines – за класическия Macintosh Desktop Interface.
    \item OSF/Motif Style Guide – за UNIX системи с Motif.
    \item OPEN LOOK Graphical Interface Functional Specification (UNIX International).
\end{itemize}

\subsection{Обработка на взаимодействията}

Добре изграденият интерфейс трябва да предлага:
\begin{itemize}
    \item \textbf{Ясна и подробна, но ненатрапчива информация} – не повече от необходимото за текущата задача. Използва се наслояване (например разкриващи се панели).
    \item \textbf{Еднозначни и лесни указания} – надписите на бутоните и връзките трябва точно да описват действието.
    \item \textbf{Смислени и информативни съобщения за грешки} – вместо „Грешка 0x8004005“ се казва „Не можете да изтриете последния администраторски акаунт. Първо създайте друг администратор.“
    \item \textbf{Лесен достъп до допълнителна информация и указания} – контекстна помощ, подсказки (tooltips), често задавани въпроси.
    \item \textbf{Възможност за отмяна на действията (Undo/Redo)} – минимално изискване за безопасност.
\end{itemize}

\subsection{Интерактивни методи за проектиране}

Интерактивните методи акцентират върху непрекъснатото включване на потребителя в процеса:
\begin{itemize}
    \item \textbf{Разбиране на нивото (опит) на потребителите} – начинаещи, средно напреднали, експерти. Интерфейсът трябва да поддържа и трите групи (напр. съкращения за експерти, стъпкови съветници за начинаещи).
    \item \textbf{Определяне на задачите} – провеждане на наблюдения, интервюта, анкети и анализ на съществуващи системи.
    \item \textbf{Избиране на стил на взаимодействие} – според контекста (например директно манипулиране за графични редактори, команден ред за системни администратори).
    \item \textbf{Освобождаване на потребителя} – системата да не принуждава потребителя да помни информация между различни стъпки (принцип на разпознаване, а не на припомняне).
    \item \textbf{Съчетаване на автоматичност и управление на потребителя} – автоматичните функции (напр. проверка на правописа) не трябва да отнемат контрола от потребителя; винаги да има възможност за ръчно отменяне или коригиране.
\end{itemize}

\subsection{Прототипиране}

Прототипът е работещ модел на интерфейса, създаден за демонстрация и оценка преди финалната реализация. Целта му е да позволи на потребителя да добие пряк опит с продукта, без да е необходимо пълно програмиране.

\textbf{Измерения на прототипите:}
\begin{itemize}
    \item \textbf{Хоризонтален прототип} – покрива много функции, но повърхностно (широчина, малка дълбочина). Подходящ за демонстриране на цялостния поток.
    \item \textbf{Вертикален прототип} – една или две функции са реализирани напълно (дълбочина), останалите са схематични. Подходящ за тестване на сложни взаимодействия.
\end{itemize}
\textbf{Степен на прецизност:}
\begin{itemize}
    \item \textbf{Ниска прецизност} – хартиени скици, схематични рисунки. Бързи и евтини, фокус върху структурата.
    \item \textbf{Висока прецизност} – компютърно генерирани екрани с реални шрифтове, цветове и интерактивност. По-близо до крайния продукт, но по-скъпи и бавни за промяна.
\end{itemize}

Процесът на прототипиране включва два основни етапа:
\begin{enumerate}
    \item Създаване на прости (напр. хартиени) прототипи – за бърза проверка на концепциите и основните работни потоци.
    \item Разработване на по-сложни автоматизирани версии – за тестване на визуалния дизайн, интеракцията и специфични детайли преди окончателното кодиране.
\end{enumerate}

Създаването на прототипи и тяхното тестване с потребители е задължително на всеки етап от проектирането. То позволява откриване на грешки и подобрения рано, когато цената на промяната е ниска.

\section*{Заключение}

Проектирането на потребителски интерфейс е сложна, мултидисциплинарна дейност, която изисква систематично прилагане на модели, методи и техники. От анализа на задачите и познавателните модели, през интерактивните стилове и визуалния дизайн на цветове, звуци и графика, до прототипирането и експерименталното оценяване – всяка стъпка трябва да бъде ориентирана към крайния потребител. Следването на описаните в конспекта принципи и използването на съвременните инструменти води до създаването на използваеми, ефективни, безопасни и приятни за работа човеко-машинни интерфейси.

\end{FlushLeft}
\end{document}
